\section*{Hoofdstuk \ref{ch:g10:drainage-basins} --- Stroomgebieden, rivieren en hoogwater}

\begin{solution}{exo:g10:drainage-basins:1}
Het \omterm{def:g10:drainage-basins:drainage-basin}{stroomgebied} van een punt van een rivier is het hele landoppervlak waarvan
het afstromende water stroomopwaarts van dat punt naar de rivier stroomt en dat
punt passeert. De \omterm{def:g10:drainage-basins:drainage-basin}{waterscheiding} is de lijn, langs de kammen, die twee
aangrenzende \omterm{def:g10:drainage-basins:drainage-basin}{stroomgebieden} scheidt. Twee druppels die een meter uit elkaar aan
weerszijden van een \omterm{def:g10:drainage-basins:drainage-basin}{waterscheiding} vallen, lopen langs tegenovergestelde
hellingen naar beneden, naar twee verschillende \omterm{def:g10:drainage-basins:drainage-basin}{stroomgebieden}: hun rivieren
kunnen verschillende zeeën bereiken.
\end{solution}

\begin{solution}{exo:g10:drainage-basins:2}
De \omterm{def:g10:drainage-basins:discharge}{afvoer} is het volume water dat elke seconde door een dwarsdoorsnede van de
rivier stroomt, in \unit{m^3/s}. \qty{1}{m^3/s} is \qty{1000}{L} per
seconde; in een dag $1000 \times 86\,400 = 86.4$ miljoen liter.
\end{solution}

\begin{solution}{exo:g10:drainage-basins:3}
Voor het hoogwater voerde de rivier ongeveer \qty{330}{m^3/s} af (tussen ongeveer
\num{240} en \qty{390}{m^3/s}). De piek was \qty{9231}{m^3/s} op
21 januari. De stijging begon op 18 januari (\qty{362}{m^3/s}): ze duurde drie
dagen.
\end{solution}

\begin{solution}{exo:g10:drainage-basins:4}
$1/50 = 0.02$, \qty{2}{\percent}. Nee: de jaren zijn onafhankelijk; de kans is
dit jaar opnieuw \qty{2}{\percent}, wat er vorig jaar ook gebeurde.
\end{solution}

\begin{solution}{exo:g10:drainage-basins:5}
Eén kubieke voet is $0.3048^3 \approx \qty{0.02832}{m^3}$, dus
$593\,000 \times 0.02832 \approx \qty{16800}{m^3/s}$. De Amazone voert
$206\,000/16\,800 \approx 12$ keer meer af.
\end{solution}

\begin{solution}{exo:g10:drainage-basins:6}
Oppervlakten van de stroken: $1.5 \times 0.3 = 0.45$, $0.90$, $1.05$ en
\qty{0.60}{m^2}: $A = \qty{3.0}{m^2}$. \omterm{def:g10:drainage-basins:discharge}{Afvoeren}: $0.2 \times 0.45 = 0.09$, $0.5 \times 0.90
= 0.45$, $0.6 \times 1.05 = 0.63$, $0.3 \times 0.60 = 0.18$: $Q =
\qty{1.35}{m^3/s}$. Gemiddelde snelheid $1.35/3.0 = \qty{0.45}{m/s}$.
\end{solution}

\begin{solution}{exo:g10:drainage-basins:7}
$66 \times 10^6 \times 0.020 = \qty{1.32e6}{m^3}$. Twee dagen zijn
$2 \times 86\,400 = \qty{172800}{s}$: $1.32 \times 10^6/172\,800 \approx
\qty{7.6}{m^3/s}$ gemiddeld.
\end{solution}

\begin{solution}{exo:g10:drainage-basins:8}
Gerangschikt: 120, 88, 65, 59, 51, 47, 42, 38, 35. Met $n = 9$ is $T = 10/m$: de
grootste (120) heeft $T = 10$ jaar, de tweede (88) 5 jaar, de derde (65)
ongeveer 3.3 jaar. De mediaan, 51, heeft rang 5 en $T = 2$ jaar: ze wordt één
jaar op twee geëvenaard of overschreden.
\end{solution}

\begin{solution}{exo:g10:drainage-basins:9}
De \omterm{def:g10:drainage-basins:hydrograph}{hydrograaf} van de stad stijgt eerder en veel hoger, en daalt snel: op steile
hellingen, daken en wegen zakt weinig water weg en bereikt het afstromende water
de beek snel en allemaal tegelijk; de \omterm{def:g10:drainage-basins:hydrograph}{vertragingstijd} is kort. De \omterm{def:g10:drainage-basins:hydrograph}{hydrograaf} van
het beboste \omterm{def:g10:drainage-basins:drainage-basin}{stroomgebied} stijgt later, tot een lagere, bredere piek: de bodem en
het bos nemen veel van de regen op en geven hem langzaam af, en het water stroomt
traag over zachte hellingen. De \omterm{def:g10:drainage-basins:hydrograph}{basisafvoer} blijft er na de bui hoger.
\end{solution}

\begin{solution}{exo:g10:drainage-basins:10}
$9231 \times 86\,400 \approx \qty{7.98e8}{m^3}$, verspreid over
$29\,900 \times 10^6 = \qty{2.99e10}{m^2}$: ongeveer \qty{0.027}{m},
\qty{27}{mm}. De \omterm{def:g10:drainage-basins:hydrograph}{basisafvoer} geeft $330 \times 86\,400/2.99 \times 10^{10}
\approx \qty{0.001}{m}$, ongeveer \qty{1}{mm} per dag: de piekdag voerde zo'n
28 keer meer af.
\end{solution}

\begin{solution}{exo:g10:drainage-basins:11}
$P = 1 - 0.9^{10} \approx 1 - 0.349 = 0.651$: ongeveer \qty{65}{\percent}. Het
komt helemaal niet met kans $0.9^{10} \approx 0.35$, ongeveer
\qty{35}{\percent}.
\end{solution}

\begin{solution}{exo:g10:drainage-basins:12}
De \omterm{def:g10:drainage-basins:hydrograph}{hydrograaf} is een driehoek van drie dagen lang en \qty{1200}{m^3/s} hoog: zijn
oppervlakte is $\tfrac12 \times 1200 \times 3 \times 86\,400 \approx
\qty{1.56e8}{m^3}$, \num{156} miljoen kubieke meter. Het deel boven
\qty{800}{m^3/s} is een kleinere driehoek met dezelfde vorm, drie keer zo
klein in hoogte ($400$ in plaats van $1200$) en in duur (één dag in plaats
van drie): $\tfrac12 \times 400 \times 86\,400 \approx \qty{1.73e7}{m^3}$,
\num{17} miljoen kubieke meter. Dat is minder dan de \num{60} miljoen die het
stuwmeer kan opslaan: ja.
\end{solution}

\begin{solution}{exo:g10:drainage-basins:13}
Tot aan haar oevers vol is de rivier $800/(100 \times 2) = \qty{4}{m}$ diep.
(1) De overstromingsvlakte voert \qty{400}{m^3/s} af over \qty{900}{m} met
\qty{0.5}{m/s}: het water is er $400/(900 \times 0.5) \approx \qty{0.9}{m}$ diep.
(2) Tussen de dijken: $1200/(100 \times 2) = \qty{6}{m}$ diep, \qty{2}{m}
boven haar oevers, meer dan twee keer de diepte op de overstromingsvlakte zonder
dijken.
\end{solution}

\begin{solution}{exo:g10:drainage-basins:14}
We willen $(1 - 1/T)^{30} = 0.5$, dus $1 - 1/T = 0.5^{1/30} \approx 0.97716$,
$1/T \approx 0.02284$ en $T \approx 44$ jaar.
\end{solution}

\begin{solution}{exo:g10:drainage-basins:15}
Elk vierkantje stuurt zijn water naar zijn laagste buur (rijen van noord naar
zuid, kolommen van west naar oost):
\begin{center}
\begin{tabular}{|c|c|c|c|}
\hline
uitstroompunt & west & west & west \\ \hline
noord & west & west & zuid \\ \hline
noord & noord & zuid & zuid \\ \hline
noord & oost & oost & uitstroompunt \\ \hline
\end{tabular}
\end{center}
Als je de pijlen volgt, wateren de vierkantjes 47, 39, 44, 37, 46 en 30 af naar
het uitstroompunt in het zuidoosten: 6 vierkantjes, \qty{0.06}{km^2}. De andere 10
vierkantjes wateren af naar het uitstroompunt in het noordwesten:
\qty{0.10}{km^2}. De \omterm{def:g10:drainage-basins:drainage-basin}{waterscheiding} loopt vanaf de oostrand tussen 50 en 47,
dan tussen 48 en 47, tussen 48 en 44, tussen 49 en 44, tussen 49 en 46, en tussen
51 en 46 naar de zuidrand.
\end{solution}

\begin{solution}{pb:g10:drainage-basins:1}
\textbf{1.} $11\,560 \times 1.609^2 \approx \qty{29900}{km^2}$.

\textbf{2.} De \omterm{def:g10:drainage-basins:drainage-basin}{waterscheiding} van het \omterm{def:g10:drainage-basins:drainage-basin}{stroomgebied} stroomopwaarts van het station.

\textbf{3.} $484\,000 \times 0.3048^3 \approx \qty{13700}{m^3/s}$.

\textbf{4.} $13\,705/3341 \approx 4.1$ keer.

\textbf{5.} $A = 400 \times 8 = \qty{3200}{m^2}$ en $v = Q/A = 13\,705/3200
\approx \qty{4.3}{m/s}$.

\textbf{6.} Ongeveer \qty{330}{m^3/s}.

\textbf{7.} \qty{9231}{m^3/s}, op 21 januari.

\textbf{8.} Drie dagen.

\textbf{9.} $9231 \times 86\,400/(29\,900 \times 10^6) \approx
\qty{0.027}{m}$: ongeveer \qty{27}{mm}.

\textbf{10.} Er bereikte nieuw water de rivier: een nieuwe regenbui op het
\omterm{def:g10:drainage-basins:drainage-basin}{stroomgebied} (of smeltende sneeuw) stuurde een tweede, kleinere golf afstromend
water.

\textbf{11.} Zeventien dagen na de piek. De stijging komt van het afstromende
water van de hellingen dicht bij de beken, dat snel en allemaal tegelijk
aankomt; op de dalende tak krijgt de rivier nog het water van de verste delen van
het \omterm{def:g10:drainage-basins:drainage-basin}{stroomgebied} en het water dat langzaam uit de bodem wegloopt.

\textbf{12.} $T = n + 1 = 96$ jaar.

\textbf{13.} $96/2 = 48$ jaar.

\textbf{14.} $96/3 = 32$ jaar.

\textbf{15.} $96/4 = 24$ jaar en $96/5 = 19.2$ jaar. Niet veel: hoe dan ook komt
deze \omterm{def:g10:drainage-basins:discharge}{afvoer} ongeveer eens in de twintig jaar terug.

\textbf{16.} Vijf jaren (1936, 1937, 1943, 1972 en 1996): $5/95 \approx
0.053$, ongeveer \qty{5}{\percent} per jaar, een \omterm{def:g10:drainage-basins:return-period}{herhalingstijd} van ongeveer
$96/5 \approx 19$ jaar.

\textbf{17.} $1/100 = 0.01$, \qty{1}{\percent}.

\textbf{18.} $1 - 0.01 = 0.99$.

\textbf{19.} $0.99^{30} \approx 0.740$.

\textbf{20.} $1 - 0.740 = 0.26$: \textbf{de kans op minstens één
honderdjaarshoogwater in 30 jaar is ongeveer \qty{26}{\percent}}, één op vier.
\end{solution}
